\section{Problem Formulation: Minimum Bisection and Photonic Mapping}
\label{sec:problem}

\subsection{Minimum Bisection / Balanced Graph Partitioning}

Given an undirected weighted graph $G=(V,E,w)$ with $|V|=n$ (even) and edge weights $w_{ij}\ge 0$, the Minimum Bisection problem seeks a partition $V = V_0 \cup V_1$ with $|V_0|=|V_1|=n/2$ minimizing the cut weight:
\begin{equation}
\label{eq:min_bisection}
C_G(x) = \sum_{(i,j)\in E} w_{ij} |x_i - x_j| \quad \text{subject to} \quad x \in \{0,1\}^n,\; \sum_{i=1}^n x_i = k = n/2.
\end{equation}
The binary variable $x_i=1$ indicates vertex $i \in V_0$ and $x_i=0$ indicates $i \in V_1$. The constraint $\sum_i x_i = k$ enforces balanced partitions.

This problem is NP-hard~\cite{garey1976complexity}. Classical exact solvers use branch-and-bound or ILP; heuristics include spectral partitioning~\cite{fiedler1973algebraic}, Kernighan-Lin~\cite{kernighan1970efficient}, and semidefinite programming relaxations~\cite{goemans1995improved}.

\subsection{Photonic Mapping: Fixed-Photon-Number Interferometry}

Consider an $n$-mode passive linear optical interferometer described by a unitary $U \in \mathrm{U}(n)$. Inject $k$ indistinguishable photons into the first $k$ modes: the input state is $|1,\dots,1,0,\dots,0\rangle$ with $k$ ones. The output state lives in the $k$-photon symmetric subspace $\mathcal{H}_k = \mathrm{Sym}^k(\mathbb{C}^n)$ of dimension
\begin{equation}
D = \dim \mathcal{H}_k = \binom{n+k-1}{k}.
\end{equation}
A basis for $\mathcal{H}_k$ is given by occupation number vectors $\mathbf{n} = (n_1,\dots,n_n)$ with $n_i \in \mathbb{N}_0$ and $\sum_i n_i = k$.

The transition amplitude from input occupation $\mathbf{s}$ to output occupation $\mathbf{t}$ is given by the bosonic permanent formula:
\begin{equation}
\label{eq:bosonic_amplitude}
\langle \mathbf{t} | \hat{\Gamma}_k(U) | \mathbf{s} \rangle = \frac{\mathrm{Per}(U_{\mathbf{t},\mathbf{s}})}{\sqrt{\prod_i s_i! \prod_j t_j!}},
\end{equation}
where $U_{\mathbf{t},\mathbf{s}}$ is the $k\times k$ submatrix of $U$ with rows repeated according to $\mathbf{t}$ and columns repeated according to $\mathbf{s}$. The probability of output occupation $\mathbf{t}$ is $p(\mathbf{t}) = |\langle \mathbf{t} | \hat{\Gamma}_k(U) | \mathbf{s} \rangle|^2$.

\subsection{Collision-Free Subspace and Balanced Partitions}

A collision-free output occupation satisfies $n_i \in \{0,1\}$ for all $i$. The collision-free subspace is
\begin{equation}
\mathcal{H}_k^{\mathrm{CF}} = \mathrm{span}\{ |\mathbf{n}\rangle : n_i \in \{0,1\},\; \sum_i n_i = k \},
\end{equation}
with dimension $d = \binom{n}{k}$. Each collision-free occupation $\mathbf{n}$ corresponds bijectively to a balanced partition: the set $\{i : n_i = 1\}$ has exactly $k$ elements.

\begin{table}[ht]
\centering
\caption{Fock space dimensions for balanced partitioning ($k=n/2$).}
\label{tab:fock_dimensions}
\begin{tabular}{cccccc}
\toprule
$n$ & $k$ & Full Fock $D=\binom{n+k-1}{k}$ & Collision-free $d=\binom{n}{k}$ & Ratio $d/D$ & $1/(d/D)$ \\
\midrule
4 & 2 & 10 & 6 & 0.600 & 1.67 \\
6 & 3 & 56 & 20 & 0.357 & 2.80 \\
8 & 4 & 330 & 70 & 0.212 & 4.71 \\
10 & 5 & 2002 & 252 & 0.126 & 7.94 \\
12 & 6 & 12376 & 924 & 0.075 & 13.4 \\
\bottomrule
\end{tabular}
\end{table}

Table~\ref{tab:fock_dimensions} shows the exponential gap between full and collision-free dimensions. For $n=12$, only $7.5\%$ of the Fock space is collision-free.

\subsection{Physical Success Probability}

Let $C^* = \min_{x: \sum x_i=k} C_G(x)$ be the optimal cut weight. The collision-free probability is
\begin{equation}
P_{\mathrm{CF}}(U) = \sum_{\mathbf{n} \in \mathrm{CF}} p(\mathbf{n}).
\end{equation}
The conditional probability of an optimal solution given collision-free detection is
\begin{equation}
P_{\mathrm{opt|CF}}(U) = \frac{1}{P_{\mathrm{CF}}(U)} \sum_{\substack{\mathbf{n} \in \mathrm{CF} \\ C_G(\mathbf{n})=C^*}} p(\mathbf{n}).
\end{equation}
The \textbf{physical success probability}---the probability per physical shot of obtaining an optimal balanced partition---is
\begin{equation}
\label{eq:p_success}
P_{\mathrm{success}}(U) = P_{\mathrm{CF}}(U) \cdot P_{\mathrm{opt|CF}}(U) = \sum_{\substack{\mathbf{n} \in \mathrm{CF} \\ C_G(\mathbf{n})=C^*}} p(\mathbf{n}).
\end{equation}
The expected number of physical shots to obtain one optimal solution is $N_{\mathrm{shots}} = 1/P_{\mathrm{success}}$.

\textbf{Critical distinction:} Prior work~\cite{wang2023bosonic} minimized the conditional energy $\mathbb{E}[C|\mathrm{CF}] = \sum_{\mathbf{n}\in\mathrm{CF}} p(\mathbf{n}) C_G(\mathbf{n}) / P_{\mathrm{CF}}$. This is \textbf{not} a valid hardware objective: if $P_{\mathrm{CF}}\to 0$, the conditional energy can be arbitrarily small while $P_{\mathrm{success}}\to 0$, meaning the device never produces a valid solution. The correct metric for hardware is $P_{\mathrm{success}}$ (or equivalently $N_{\mathrm{shots}}$).
